\section*{Solutions: Complex Numbers and Complex Functions Quiz}

\begin{enumerate}
    \item \textbf{Polar Form, Powers, and Roots}

    Given
    \[
        z=-8+8\sqrt{3}\,i.
    \]

    First compute the modulus:
    \[
        r=\sqrt{(-8)^2+(8\sqrt{3})^2}
        =
        \sqrt{64+192}
        =
        \sqrt{256}
        =
        16.
    \]

    Since $x=-8<0$ and $y=8\sqrt{3}>0$, the point lies in Quadrant II. Also,
    \[
        \tan\theta=\frac{8\sqrt{3}}{-8}=-\sqrt{3}.
    \]
    The reference angle is $\pi/3$, so the principal argument is
    \[
        \theta=\frac{2\pi}{3}.
    \]

    Hence,
    \[
        z=16e^{i2\pi/3}.
    \]

    We now solve
    \[
        w^4=16e^{i2\pi/3}.
    \]

    Write
    \[
        w=\rho e^{i\phi}.
    \]
    Then
    \[
        w^4=\rho^4e^{i4\phi}.
    \]

    Therefore,
    \[
        \rho^4=16,
        \qquad
        4\phi=\frac{2\pi}{3}+2k\pi.
    \]

    Thus,
    \[
        \rho=2,
    \]
    and
    \[
        \phi_k=\frac{1}{4}\left(\frac{2\pi}{3}+2k\pi\right)
        =
        \frac{\pi}{6}+\frac{k\pi}{2},
        \qquad k=0,1,2,3.
    \]

    Therefore, the four roots are
    \[
        w_k=2e^{i\left(\frac{\pi}{6}+\frac{k\pi}{2}\right)},
        \qquad k=0,1,2,3.
    \]

    Explicitly:
    \[
        w_0=2e^{i\pi/6},
        \qquad
        w_1=2e^{i2\pi/3},
    \]
    \[
        w_2=2e^{i7\pi/6},
        \qquad
        w_3=2e^{i5\pi/3}.
    \]

    Geometrically, the roots lie on the circle
    \[
        |w|=2
    \]
    and are equally spaced by
    \[
        \frac{2\pi}{4}=\frac{\pi}{2}.
    \]
    Hence, they form the vertices of a square centered at the origin.

    \item \textbf{Complex Mapping}

    We are given
    \[
        w=\frac{1}{z-i}
    \]
    and
    \[
        z=x+2i.
    \]

    Substitute into the mapping:
    \[
        w=\frac{1}{x+2i-i}
        =
        \frac{1}{x+i}.
    \]

    Rationalize:
    \[
        w=\frac{1}{x+i}\cdot \frac{x-i}{x-i}
        =
        \frac{x-i}{x^2+1}.
    \]

    Write
    \[
        w=u+iv.
    \]
    Then
    \[
        u=\frac{x}{x^2+1},
        \qquad
        v=-\frac{1}{x^2+1}.
    \]

    We eliminate $x$. Divide $u$ by $v$:
    \[
        \frac{u}{v}
        =
        \frac{x/(x^2+1)}{-1/(x^2+1)}
        =
        -x.
    \]
    Hence,
    \[
        x=-\frac{u}{v}.
    \]

    Since
    \[
        v=-\frac{1}{x^2+1},
    \]
    substitute $x=-u/v$:
    \[
        v=-\frac{1}{(u^2/v^2)+1}.
    \]

    Simplify:
    \[
        v=-\frac{v^2}{u^2+v^2}.
    \]

    Since $v\neq 0$ on the image curve, divide by $v$:
    \[
        1=-\frac{v}{u^2+v^2}.
    \]

    Therefore,
    \[
        u^2+v^2=-v.
    \]

    Rearranging:
    \[
        u^2+v^2+v=0.
    \]

    Complete the square:
    \[
        u^2+\left(v+\frac12\right)^2=\frac14.
    \]

    Hence, the image is the circle
    \[
        u^2+\left(v+\frac12\right)^2=\left(\frac12\right)^2.
    \]

    This is a circle of radius $1/2$ centered at
    \[
        0-\frac12 i.
    \]

    \item \textbf{Analyticity and the Cauchy--Riemann Equations}

    We are given
    \[
        u(x,y)=x^3-3xy^2+2x,
        \qquad
        v(x,y)=3x^2y-y^3+ay.
    \]

    Compute the partial derivatives:
    \[
        u_x=3x^2-3y^2+2,
        \qquad
        u_y=-6xy,
    \]
    \[
        v_x=6xy,
        \qquad
        v_y=3x^2-3y^2+a.
    \]

    The Cauchy--Riemann equations require
    \[
        u_x=v_y,
        \qquad
        u_y=-v_x.
    \]

    The second equation gives
    \[
        -6xy=-6xy,
    \]
    which is automatically satisfied.

    The first equation gives
    \[
        3x^2-3y^2+2=3x^2-3y^2+a.
    \]

    Therefore,
    \[
        a=2.
    \]

    With $a=2$,
    \[
        v(x,y)=3x^2y-y^3+2y.
    \]

    Now observe that
    \[
        z^3=(x+iy)^3
        =
        x^3+3x^2iy+3x(iy)^2+(iy)^3.
    \]

    Since $i^2=-1$ and $i^3=-i$,
    \[
        z^3=x^3+3ix^2y-3xy^2-iy^3.
    \]

    Hence,
    \[
        z^3=(x^3-3xy^2)+i(3x^2y-y^3).
    \]

    Also,
    \[
        2z=2x+i2y.
    \]

    Therefore,
    \[
        z^3+2z
        =
        (x^3-3xy^2+2x)
        +
        i(3x^2y-y^3+2y).
    \]

    Thus,
    \[
        f(z)=z^3+2z.
    \]

    \item \textbf{Complex Integration}

    We evaluate
    \[
        \int_C \overline{z}\,dz.
    \]

    \textbf{Path $C_1$:} Let
    \[
        z(t)=(1+i)t,
        \qquad 0\le t\le 1.
    \]

    Then
    \[
        \overline{z(t)}=(1-i)t,
        \qquad
        dz=(1+i)\,dt.
    \]

    Thus,
    \[
        \int_{C_1}\overline{z}\,dz
        =
        \int_0^1 (1-i)t(1+i)\,dt.
    \]

    Since
    \[
        (1-i)(1+i)=1-i^2=2,
    \]
    we get
    \[
        \int_{C_1}\overline{z}\,dz
        =
        \int_0^1 2t\,dt
        =
        \left[t^2\right]_0^1
        =
        1.
    \]

    \textbf{Path $C_2$:} This path has two segments.

    First segment: from $0$ to $1$,
    \[
        z(t)=t,
        \qquad 0\le t\le 1.
    \]
    Then
    \[
        \overline{z}=t,
        \qquad
        dz=dt.
    \]
    Thus,
    \[
        \int_{0}^{1} \overline{z}\,dz
        =
        \int_0^1 t\,dt
        =
        \frac12.
    \]

    Second segment: from $1$ to $1+i$,
    \[
        z(t)=1+it,
        \qquad 0\le t\le 1.
    \]
    Then
    \[
        \overline{z}=1-it,
        \qquad
        dz=i\,dt.
    \]

    Therefore,
    \[
        \int (1-it)i\,dt
        =
        \int (i+t)\,dt.
    \]

    Hence,
    \[
        \int_0^1 (i+t)\,dt
        =
        i+\frac12.
    \]

    Adding both segments:
    \[
        \int_{C_2}\overline{z}\,dz
        =
        \frac12+\left(i+\frac12\right)
        =
        1+i.
    \]

    Since
    \[
        \int_{C_1}\overline{z}\,dz=1
    \]
    but
    \[
        \int_{C_2}\overline{z}\,dz=1+i,
    \]
    the integral depends on the path.

    Therefore, $\overline{z}$ does not have path-independent integrals.

    \item \textbf{Complex Exponential and Logarithm}

    We solve
    \[
        e^z=-2+2i.
    \]

    First express the right-hand side in polar form:
    \[
        -2+2i=2\sqrt{2}\,e^{i3\pi/4}.
    \]

    Since the argument is multivalued,
    \[
        -2+2i=2\sqrt{2}\,e^{i(3\pi/4+2k\pi)},
        \qquad k\in\mathbb{Z}.
    \]

    Let
    \[
        z=x+iy.
    \]
    Then
    \[
        e^z=e^{x+iy}=e^xe^{iy}.
    \]

    Matching modulus and argument gives
    \[
        e^x=2\sqrt{2},
        \qquad
        y=\frac{3\pi}{4}+2k\pi.
    \]

    Therefore,
    \[
        x=\ln(2\sqrt{2}).
    \]

    Hence all solutions are
    \[
        z=\ln(2\sqrt{2})+i\left(\frac{3\pi}{4}+2k\pi\right),
        \qquad k\in\mathbb{Z}.
    \]

    Since
    \[
        \ln(2\sqrt{2})
        =
        \ln(2^{3/2})
        =
        \frac32\ln 2,
    \]
    we may also write
    \[
        z=\frac32\ln 2+i\left(\frac{3\pi}{4}+2k\pi\right),
        \qquad k\in\mathbb{Z}.
    \]

    The principal solution uses
    \[
        \operatorname{Arg}(-2+2i)=\frac{3\pi}{4}.
    \]

    Therefore,
    \[
        z_{\text{principal}}=\frac32\ln 2+i\frac{3\pi}{4}.
    \]
\end{enumerate}
